% Ad Familiares 8.12 (ad-familiares-08-12) --- English translation
% Translated by Claude (Anthropic) from the Latin (Perseus).
% Style guide: STYLE.md.

\ciceroLetterOpener{CAELIVS CICERONI S.}

\ciceroSection{1} It shames me to confess to you, and to complain to you, of the wrongs of Appius --- a most ungrateful man, who began to hate me because he owed me great kindnesses, and who, when as a man of avarice he could not bring himself to pay them, declared a secret war on me --- though secret in such a way that many men reported it to me, and I could easily see for myself that he was thinking ill of me. But when I learned for certain that he had been sounding out the augural college, and then was speaking openly with certain men, was holding counsel with L. Domitius (a man, as things stand, my bitterest enemy), and meant to make a present of this little gift to Cn. Pompey, then I went to find him, and to beg him off from doing me an injury, when I had thought him to owe me his very life --- and I could not get it out of him.

\ciceroSection{2} What then? Even so I spoke to a number of friends, the witnesses to my services to him. After I had perceived that the man did not even hold me worthy of his giving satisfaction to anyone, I preferred to put myself under obligation to his colleague --- a man as remote from me as could be, and, on account of his friendship with you, not the best disposed to me --- rather than submit to the looks of that ape. When he found that out, he flared up and shouted that I was casting about for a ground of quarrel, so that, if he were to fail of satisfying me in the money matter, I should pursue him under this colour of feud. After that he was forever sending for Pola Servius to bring the prosecution, and laying his plots with Domitius in earnest.

\ciceroSection{3} When they were getting nowhere in their attempt to set a prosecutor on me under any other law, they thought to indict me under the very one under which they had no case to bring: with monumental insolence the men saw to it that, at the height of my games at the Circus, I should be summoned under the Lex Scantinia (a statute against unnatural offences). Pola had scarcely uttered the indictment when I served Appius as censor with summons under the same law. Nothing could have fallen out better: for it has been so applauded by the people --- and not by the lowest sort either --- that the rumour of it has stung Appius worse than the summons itself. Besides, I have begun to demand of him the chapel that stands in his house.

\ciceroSection{4} The delay of this slave of mine who brought the letter to you puts me out: after he had received the earlier letters, he stayed on more than forty days. What I am to write to you I do not know. You know that Domitius is in dread of the day of the elections. I look for you eagerly and long to see you at the earliest. From you I ask this --- that you grieve for my wrongs as you reckon I am wont both to grieve for yours and to avenge them.
